\chapter{Test e Risultati}

In questo capitolo vengono presentati i test eseguiti sulla struttura 
\texttt{HashRBTree} per verificare la correttezza delle operazioni di 
inserimento, ricerca, cancellazione e stampa.

\section{Ambiente di Test}

I test sono stati eseguiti nel seguente ambiente:

\begin{itemize}
    \item \textbf{Sistema operativo}: Linux Mint
    \item \textbf{Compilatore}: GCC 13.3.0
    \item \textbf{Standard}: C++17
    \item \textbf{Build system}: CMake 3.10
\end{itemize}

\section{Test 1 --- Inserimento e Stampa}

Vengono inserite sette coppie chiave-valore e viene verificata la struttura 
dell'albero tramite \texttt{print()}.

\begin{lstlisting}
tree.insert("apple",      "frutto rosso");
tree.insert("banana",     "frutto giallo");
tree.insert("cherry",     "frutto rosso piccolo");
tree.insert("date",       "frutto dolce");
tree.insert("elderberry", "frutto viola");
tree.insert("fig",        "frutto mediterraneo");
tree.insert("grape",      "frutto a grappolo");
\end{lstlisting}

\textbf{Output ottenuto}:

\begin{lstlisting}
+-- [1938937894 | B] banana:frutto giallo
    +-- [253337143 | R] apple:frutto rosso
    |   +-- [193491643 | B] fig:frutto mediterraneo
    |   `-- [735886645 | B] elderberry:frutto viola
    |       +-- [260508340 | R] grape:frutto a grappolo
    `-- [1986069970 | B] cherry:frutto rosso piccolo
        `-- [2090176867 | R] date:frutto dolce
\end{lstlisting}

\textbf{Verifica}: l'albero rispetta tutte le proprietà del Red-Black Tree:
\begin{itemize}
    \item la radice è \texttt{BLACK};
    \item non esistono due nodi \texttt{RED} consecutivi;
    \item il black-height è uniforme su tutti i percorsi dalla radice 
    alle foglie.
\end{itemize}

\section{Test 2 --- Ricerca}

\begin{lstlisting}
tree.search("apple",  "");
tree.search("banana", "");
tree.search("mango",  "");
\end{lstlisting}

\textbf{Output ottenuto}:

\begin{lstlisting}
Trovato: apple -> frutto rosso
Trovato: banana -> frutto giallo
Chiave "mango" non trovata.
\end{lstlisting}

\textbf{Verifica}: le chiavi presenti vengono trovate correttamente; 
la ricerca di una chiave assente restituisce il messaggio di errore 
senza crash.

\section{Test 3 --- Gestione delle Collisioni}

Vengono inserite due chiavi aggiuntive per verificare il comportamento 
in caso di collisioni hash.

\begin{lstlisting}
tree.insert("listen", "ascoltare");
tree.insert("silent", "silenzioso");
\end{lstlisting}

\textbf{Verifica}: la funzione djb2 è sensibile all'ordine dei caratteri, 
quindi \texttt{listen} e \texttt{silent} producono hash distinti e vengono 
inseriti come nodi separati. L'albero si ribilancia correttamente dopo 
entrambi gli inserimenti.

\section{Test 4 --- Cancellazione}

\begin{lstlisting}
tree.remove("banana", "");
tree.remove("mango",  "");
\end{lstlisting}

\textbf{Output ottenuto}:

\begin{lstlisting}
Chiave "mango" non trovata.

+-- [1986069970 | B] cherry:frutto rosso piccolo
    +-- [253337143 | R] apple:frutto rosso
    |   +-- [193491643 | B] fig:frutto mediterraneo
    |   |   +-- [192495636 | R] listen:ascoltare
    |   `-- [466175796 | B] silent:silenzioso
    |       +-- [260508340 | R] grape:frutto a grappolo
    |       `-- [735886645 | R] elderberry:frutto viola
    `-- [2090176867 | B] date:frutto dolce
\end{lstlisting}

\textbf{Verifica}: la rimozione di \texttt{banana} (nodo radice) viene 
gestita correttamente tramite il successore in-order; l'albero rimane 
bilanciato. La rimozione di \texttt{mango} (chiave assente) produce 
il messaggio di errore senza alterare la struttura.

\section{Test 5 --- Cancellazioni Multiple}

\begin{lstlisting}
tree.remove("apple",  "");
tree.remove("cherry", "");
\end{lstlisting}

\textbf{Output ottenuto}:

\begin{lstlisting}
+-- [260508340 | B] grape:frutto a grappolo
    +-- [193491643 | B] fig:frutto mediterraneo
    |   +-- [192495636 | R] listen:ascoltare
    `-- [735886645 | R] elderberry:frutto viola
        +-- [466175796 | B] silent:silenzioso
        `-- [2090176867 | B] date:frutto dolce
\end{lstlisting}

\textbf{Verifica}: dopo tre cancellazioni consecutive l'albero rimane 
correttamente bilanciato, confermando il corretto funzionamento di 
\texttt{deleteFix} in scenari di stress.

\section{Considerazioni Finali}

I test confermano che la struttura \texttt{HashRBTree} si comporta 
correttamente in tutti gli scenari verificati:

\begin{itemize}
    \item le proprietà del Red-Black Tree vengono mantenute dopo ogni 
    inserimento e cancellazione;
    \item la gestione delle collisioni tramite lista funziona correttamente;
    \item i casi limite (chiave assente, albero con un solo nodo, 
    cancellazione della radice) sono gestiti senza errori.
\end{itemize}